\subsection{Duality}

\remark{For the rest of the solutions manual, unless otherwise stated, we will be using the Kronecker delta notation, i.e., we define $\delta_{ij} = 1$ if $i = j$ and $\delta_{ij} = 0$ if $i \neq j$.}

\begin{exercise}[3f1]
    Explain why each linear functional is surjective or is the zero map.
\end{exercise}

\begin{sol}
    Let $\varphi \in V'$. If $\varphi = 0$, then it is the zero map. Otherwise, there exists $v \in V$ such that $\varphi(v) \neq 0$, say $\varphi(v) = c \neq 0$. Then for any $x \in \f$, we have
    \[\varphi\left(\frac{x}{c} v\right) = \frac{x}{c} \varphi(v) = \frac{x}{c} c = x.\]
    Hence, $\varphi$ is surjective.
\end{sol}

\begin{exercise}[3f2]
    Give three distinct examples of linear functionals on $\r^{[0, 1]}$.
\end{exercise}

\begin{sol}
    Consider the linear functionals on $\r^{[0, 1]}$ defined by
    \[\varphi_1(f) = f(0), \quad \varphi_2(f) = f(1), \quad \varphi_3(f) = f\left(\frac 12\right). \qh\]
\end{sol}

\begin{exercise}[3f3]
    Suppose $V$ is finite-dimensional and $v \in V$ with $v \neq 0$. Prove that there exists $\varphi \in V'$ such that $\varphi(v) = 1$.
\end{exercise}

\begin{sol}
    Since $V$ is finite-dimensional, we can extend $v$ to a basis of $V$, say $v, v_2, \ldots, v_n$. Define $\varphi \in V'$ by setting $\varphi(v) = 1$ and $\varphi(v_i) = 0$ for $i = 2, \ldots, n$. This defines a linear functional because it is defined on a basis and extended linearly. Clearly, $\varphi(v) = 1$.
\end{sol}

\begin{exercise}[3f4]
    Suppose $V$ is finite-dimensional and $U$ is a subspace of $V$ such that $U \neq V$. Prove that there exists $\varphi \in V'$ such that $\varphi(u) = 0$ for every $u \in U$ but $\varphi \neq 0$.
\end{exercise}

\begin{sol}
    Since $U \neq V$, there exists $v_1 \in V$ such that $v_1 \notin U$. Extend a basis of $U$ to a basis of $V$, say $u_1, \ldots, u_k, v_1, \ldots, v_m$. Define $\varphi \in V'$ by setting $\varphi(u_i) = 0$ for all $i = 1, \ldots, k$ and $\varphi(v_j) = 0$ for all $j = 2, \ldots, m$ but $\varphi(v_1) = 1$. This defines a linear functional because it is defined on a basis and extended linearly. Clearly, $\varphi(u) = 0$ for all $u \in U$ and $\varphi \neq 0$.
\end{sol}

\begin{exercise}[3f5]
    Suppose $T \in \lin{V, W}$ and $w_1, \ldots, w_m$ is a basis of $\range T$. Hence for each $v \in V$, there exist unique numbers $\varphi_1(v), \ldots, \varphi_m(v)$ such that
    \[Tv = \varphi_1(v) w_1 + \cdots + \varphi_m(v) w_m,\]
    thus defining functions $\varphi_1, \ldots, \varphi_m$ from $V$ to $\f$. Show that each of the functions $\varphi_1, \ldots, \varphi_m$ is a linear functional on $V$.
\end{exercise}

\begin{sol}
    Let $v \in V$. Since $w_1, \ldots, w_m$ is a basis of $\range T$, there exist unique scalars $\varphi_1(v), \ldots, \varphi_m(v)$ such that
    \[Tv = \varphi_1(v) w_1 + \cdots + \varphi_m(v) w_m.\]
    To show that each $\varphi_i$ is a linear functional, let $v, u \in V$ and $a, b \in \f$. Then
    \[T(av + bu) = aTv + bTu = a(\varphi_1(v) w_1 + \cdots + \varphi_m(v) w_m) + b(\varphi_1(u) w_1 + \cdots + \varphi_m(u) w_m).\]
    By the uniqueness of the representation in terms of the basis $w_1, \ldots, w_m$, we must have
    \[\varphi_i(av + bu) = a\varphi_i(v) + b\varphi_i(u) \quad \text{for each } i = 1, \ldots, m.\]
    Hence, each $\varphi_i$ is a linear functional on $V$.
\end{sol}

\begin{exercise}[3f6]
    Suppose $\varphi, \beta \in V'$. Prove that $\nul \varphi \subseteq \nul \beta$ if and only if there exists $c \in \f$ such that $\beta = c\varphi$.
\end{exercise}

\begin{sol}
    Suppose $\nul \varphi \subseteq \nul \beta$. If $\varphi = 0$, then $\nul \varphi = V$, and hence $\nul \beta = V$, which implies $\beta = 0$. In this case, we can take $c = 0$ and have $\beta = c\varphi$. Now assume $\varphi \neq 0$. Pick $v_0 \in V$ such that $\varphi(v_0) \neq 0$. Since $v_0 \notin \nul \varphi$, we have $v_0 \notin \nul \beta$ unless $\beta = 0$. Define
    \[c = \frac{\beta(v_0)}{\varphi(v_0)}.\] 
    For any $v \in V$, consider
    \[v - \frac{\varphi(v)}{\varphi(v_0)} v_0.\]
    This vector is in $\nul \varphi$, so it is also in $\nul \beta$. Hence,
    \[\beta\left(v - \frac{\varphi(v)}{\varphi(v_0)} v_0\right) = 0 \implies \beta(v) = \frac{\varphi(v)}{\varphi(v_0)} \beta(v_0) = c \varphi(v).\]
    Therefore, $\beta = c\varphi$.
    
    Conversely, if $\beta = c\varphi$, then clearly $\nul \varphi \subseteq \nul \beta$ because if $\varphi(v) = 0$, then $\beta(v) = c\varphi(v) = 0$.
\end{sol}

\begin{exercise}[3f7]
    Suppose that $V_1, \ldots, V_m$ are vector spaces. Prove that $(V_1 \times \cdots \times V_m)'$ and $V_1' \times \cdots \times V_m'$ are isomorphic vector spaces.
\end{exercise}

\begin{sol}
    Define a map $\Phi: (V_1 \times \cdots \times V_m)' \to V_1' \times \cdots \times V_m'$ by
    \[\Phi(\varphi) = (\varphi_1, \ldots, \varphi_m),\]
    where $\varphi_i \in V_i'$ is defined by $\varphi_i(v_i) = \varphi(0, \ldots, 0, v_i, 0, \ldots, 0)$, with $v_i$ in the $i$th position. 

    First, we show that $\Phi$ is linear. Let $\varphi, \psi \in (V_1 \times \cdots \times V_m)'$ and $a, b \in \f$. Then
    \[\Phi(a\varphi + b\psi) = ((a\varphi + b\psi)_1, \ldots, (a\varphi + b\psi)_m) = (a\varphi_1 + b\psi_1, \ldots, a\varphi_m + b\psi_m) = a\Phi(\varphi) + b\Phi(\psi).\]

    Next, we show that $\Phi$ is bijective. For injectivity, suppose $\Phi(\varphi) = 0$. Then $\varphi_i = 0$ for all $i$, which implies $\varphi(0, \ldots, 0, v_i, 0, \ldots, 0) = 0$ for all $v_i \in V_i$. By linearity, $\varphi$ must be zero on all of $V_1 \times \cdots \times V_m$, so $\varphi = 0$. For surjectivity, given $(\varphi_1, \ldots, \varphi_m) \in V_1' \times \cdots \times V_m'$, define $\varphi \in (V_1 \times \cdots \times V_m)'$ by
    \[\varphi(v_1, \ldots, v_m) = \varphi_1(v_1) + \cdots + \varphi_m(v_m).\]
    Then $\Phi(\varphi) = (\varphi_1, \ldots, \varphi_m)$, so $\Phi$ is surjective.

    Therefore, $\Phi$ is a linear isomorphism, and $(V_1 \times \cdots \times V_m)'$ and $V_1' \times \cdots \times V_m'$ are isomorphic vector spaces.
\end{sol}

\begin{exercise}[3f8]
    Suppose $v_1, \ldots, v_n$ is a basis of $V$ and $\varphi_1, \ldots, \varphi_n$ is the dual basis of $V'$. Define $\Gamma: V \to \f^n$ and $\Lambda: \f^n \to V$ by
    \[\Gamma(v) = (\varphi_1(v), \ldots, \varphi_n(v)) \quad \text{and} \quad \Lambda(a_1, \ldots, a_n) = a_1 v_1 + \cdots + a_n v_n.\]
    Explain why $\Gamma$ and $\Lambda$ are inverses of each other.
\end{exercise}

\begin{sol}
    To show that $\Gamma$ and $\Lambda$ are inverses of each other, we need to check that $\Gamma(\Lambda(a_1, \ldots, a_n)) = (a_1, \ldots, a_n)$ for all $(a_1, \ldots, a_n) \in \f^n$ and $\Lambda(\Gamma(v)) = v$ for all $v \in V$.

    First, let $(a_1, \ldots, a_n) \in \f^n$. Then
    \[\Gamma(\Lambda(a_1, \ldots, a_n)) = \Gamma(a_1 v_1 + \cdots + a_n v_n) = (\varphi_1(a_1 v_1 + \cdots + a_n v_n), \ldots, \varphi_n(a_1 v_1 + \cdots + a_n v_n)).\]
    By the definition of the dual basis, we have $\varphi_i(v_j) = \delta_{ij}$ for all $i, j$. Therefore,
    \[\varphi_i(a_1 v_1 + \cdots + a_n v_n) = a_i.\]
    Therefore,
    \[\Gamma(\Lambda(a_1, \ldots, a_n)) = (a_1, \ldots, a_n).\]
    Next, let $v \in V$. Write $v = b_1 v_1 + \cdots + b_n v_n$ for some scalars $b_1, \ldots, b_n$. Then
    \[\Lambda(\Gamma(v)) = \Lambda(\varphi_1(v), \ldots, \varphi_n(v)) = \Lambda(b_1, \ldots, b_n) = b_1 v_1 + \cdots + b_n v_n = v.\]

    Hence, $\Gamma$ and $\Lambda$ are inverses of each other.
\end{sol}

\begin{exercise}[3f9]
    Suppose $m$ is a positive integer. Show that the dual basis of the basis $1, x, \ldots, x^m$ of $\pp_m(\r)$ is $\varphi_0, \varphi_1, \ldots, \varphi_m$, where
    \[\varphi_k(p) = \frac{p^{(k)}(0)}{k!}.\]
    
    \note{Here $p^{(k)}$ denotes the $k$th derivative of $p$, with the understanding that the $0$th derivative of $p$ is $p$.}
\end{exercise}

\begin{sol}
    Let $p \in \pp_m(\r)$. We can write $p(x) = a_0 + a_1 x + \cdots + a_m x^m$ for some scalars $a_0, \ldots, a_m$. Then for each $k = 0, 1, \ldots, m$, we have
    \begin{align*}
        p(x) &= a_0 + a_1 x + a_2 x^2 + \cdots + a_m x^m \\
        p'(x) &= a_1 + 2a_2 x + \cdots + ma_m x^{m - 1} \\
        p''(x) &= 2a_2 + \cdots + m(m - 1)a_m x^{m - 2} \\
        &\;\vdots \\
        p^{(k)}(x) &= k!a_k + \cdots + \frac{m!}{(m - k)!} a_m x^{m - k}
    \end{align*}
    Then
    \[\varphi_k(p) = \frac{p^{(k)}(0)}{k!} = a_k.\]
    Therefore, $\varphi_k(x^j) = \delta_{kj}$ for all $0 \leq k, j \leq m$, which shows that $\varphi_0, \varphi_1, \ldots, \varphi_m$ is indeed the dual basis of $1, x, \ldots, x^m$.
\end{sol}

\begin{exercise}[3f10]
    Suppose $m$ is a positive integer.
    \begin{subexercises}
        \item Show that $1, x - 5, \ldots, (x - 5)^m$ is a basis of $\pp_m(\r)$.
        \item What is the dual basis of the basis in (a)?
    \end{subexercises}
\end{exercise}

\begin{sole}
    \item Note that $1, x - 5, \ldots, (x - 5)^m$ is a list of $m + 1$ polynomials in $\pp_m(\r)$ with differing degrees, and since $\pp_m(\r)$ has dimension $m + 1$, this list forms a basis of $\pp_m(\r)$.
    \item Substitute $y = x - 5$. Then the basis $1, x - 5, \ldots, (x - 5)^m$ becomes $1, y, \ldots, y^m$. By \autoref{ex:3f9}, the dual basis of $1, y, \ldots, y^m$ is $\varphi_0, \varphi_1, \ldots, \varphi_m$, where
    \[\varphi_k(p) = \frac{p^{(k)}(0)}{k!}.\]
    Since $y = x - 5$, then
    \[\varphi_k(p) = \frac{p^{(k)}(5)}{k!}. \qh\]
\end{sole}

\begin{exercise}[3f11]
    Suppose $v_1, \ldots, v_n$ is a basis of $V$ and $\varphi_1, \ldots, \varphi_n$ is the corresponding dual basis of $V'$. Suppose $\psi \in V'$. Prove that
    \[\psi = \psi(v_1) \varphi_1 + \cdots + \psi(v_n) \varphi_n.\]
\end{exercise}

\begin{sol}
    Let $v \in V$. Since $\varphi_1, \ldots, \varphi_n$ is the dual basis of $v_1, \ldots, v_n$, we have $\varphi_i(v_j) = \delta_{ij}$ for all $1 \leq i, j \leq n$. Let $v = \alpha_1 v_1 + \cdots + \alpha_n v_n$ for some scalars $\alpha_1, \ldots, \alpha_n$. Then
    \begin{align*}
        (\psi(v_1) \varphi_1 + \cdots + \psi(v_n) \varphi_n)(v) &= \psi(v_1) \varphi_1(v) + \cdots + \psi(v_n) \varphi_n(v) \\
        &= \psi(v_1) \alpha_1 + \cdots + \psi(v_n) \alpha_n \\
        &= \psi(\alpha_1 v_1 + \cdots + \alpha_n v_n) \\
        &= \psi(v)
    \end{align*}
    Since this holds for all $v \in V$, we conclude that
    \[\psi = \psi(v_1) \varphi_1 + \cdots + \psi(v_n) \varphi_n. \qh\]
\end{sol}

\begin{exercise}[3f12]
    Suppose $S, T \in \lin{V, W}$.
    \begin{subexercises}
        \item Prove that $(S + T)' = S' + T'$.
        \item Prove that $(\lambda T)' = \lambda T'$ for all $\lambda \in \f$.
    \end{subexercises}
\end{exercise}

\begin{sole}
    \item Let $\psi \in W'$. Then for any $v \in V$, we have
    \begin{align*}
        ((S + T)'(\psi))(v) &= \psi((S + T)(v)) \\
        &= \psi(S(v) + T(v)) \\
        &= \psi(S(v)) + \psi(T(v)) \\
        &= (S'(\psi))(v) + (T'(\psi))(v) \\
        &= (S'(\psi) + T'(\psi))(v).
    \end{align*}
    Since this holds for all $v \in V$, we conclude that
    \[(S + T)'(\psi) = S'(\psi) + T'(\psi).\]
    Since this holds for all $\psi \in W'$, we have
    \[(S + T)' = S' + T'.\]
    \item Let $\psi \in W'$. Then for any $v \in V$, we have
    \begin{align*}
        ((\lambda T)'(\psi))(v) &= \psi((\lambda T)(v)) \\
        &= \psi(\lambda T(v)) \\
        &= \lambda \psi(T(v)) \\
        &= (\lambda T'(\psi))(v).
    \end{align*}
    Since this holds for all $v \in V$, we conclude that
    \[(\lambda T)'(\psi) = \lambda T'(\psi).\]
    Since this holds for all $\psi \in W'$, we have
    \[(\lambda T)' = \lambda T'. \qh\]
\end{sole}

\begin{exercise}[3f13]
    Show that the dual map of the identity operator on $V$ is the identity operator on $V'$.
\end{exercise}

\begin{sol}
    Let $\psi \in V'$. Let $I'$ denote the dual map of the identity operator $I$ on $V$. Then for any $v \in V$, we have
    \begin{align*}
        (I'(\psi))(v) &= \psi(I(v)) \\
        &= \psi(v).
    \end{align*}
    Since this holds for all $v \in V$, we conclude that
    \[I'(\psi) = \psi.\]
    Since this holds for all $\psi \in V'$, then $I'$ must be the identity operator on $V'$.
\end{sol}

\begin{exercise}[3f14]
    Define $T: \r^3 \to \r^2$ by
    \[T(x, y, z) = (4x + 5y + 6z, 7x + 8y + 9z).\]
    Suppose $\varphi_1, \varphi_2$ denotes the dual basis of the standard basis of $\r^2$ and $\psi_1, \psi_2, \psi_3$ denotes the dual basis of the standard basis of $\r^3$.
    \begin{subexercises}
        \item Describe the linear functionals $T'(\varphi_1)$ and $T'(\varphi_2)$.
        \item Write $T'(\varphi_1)$ and $T'(\varphi_2)$ as linear combinations of $\psi_1, \psi_2, \psi_3$.
    \end{subexercises}
\end{exercise}

\begin{sole}
    \item Let $(x, y, z) \in \r^3$. Then we have
    \begin{align*}
        (T'(\varphi_1))((x, y, z)) &= \varphi_1(T(x, y, z)) \\
        &= \varphi_1(4x + 5y + 6z, 7x + 8y + 9z) \\
        &= 4x + 5y + 6z, \\
        (T'(\varphi_2))((x, y, z)) &= \varphi_2(T(x, y, z)) \\
        &= \varphi_2(4x + 5y + 6z, 7x + 8y + 9z) \\
        &= 7x + 8y + 9z.
    \end{align*}
    \item We can write $T'(\varphi_1)$ and $T'(\varphi_2)$ as linear combinations of $\psi_1, \psi_2, \psi_3$ as follows:
    \begin{align*}
        T'(\varphi_1) &= 4\psi_1 + 5\psi_2 + 6\psi_3, \\
        T'(\varphi_2) &= 7\psi_1 + 8\psi_2 + 9\psi_3.
    \end{align*} \qh
\end{sole}

\begin{exercise}[3f15]
    Define $T: \pp(\r) \to \pp(\r)$ by
    \[(Tp)(x) = x^2 p(x) + p''(x)\]
    for each $x \in \r$.
    \begin{subexercises}
        \item Suppose $\varphi \in \pp(\r)'$ is defined by $\varphi(p) = p'(4)$. Describe the linear functional $T'(\varphi)$ on $\pp(\r)$.
        \item Suppose $\varphi \in \pp(\r)'$ is defined by
        \[\varphi(p) = \int_0^1 p.\]
        Evaluate $(T'(\varphi))(x^3)$.
    \end{subexercises}
\end{exercise}

\begin{sole}
    \item Let $p \in \pp(\r)$. Then we have
    \begin{align*}
        (T'(\varphi))(p) &= \varphi(T(p)) \\
        &= (Tp)'(4) \\
        &= (2x p(x) + x^2 p'(x) + p'''(x))\big|_{x=4}. \\
        &= 8 p(4) + 16 p'(4) + p'''(4).
    \end{align*}
    \item Let $p(x) = x^3$. Then we have
    \begin{align*}
        (T'(\varphi))(x^3) &= \varphi(T(x^3)) \\
        &= \int_0^1 T(x^3) \dd x \\
        &= \int_0^1 (x^2 \cdot x^3 + (x^3)'') \dd x \\
        &= \int_0^1 (x^5 + 6x) \dd x\\
        &= \frac{1}{6} + \frac{18}{6} \\
        &= \frac{19}{6}. \qh
    \end{align*}
\end{sole}

\begin{exercise}[3f16]
    Suppose $W$ is finite-dimensional and $T \in \lin{V, W}$. Prove that
    \[T' = 0 \iff T = 0.\]
\end{exercise}

\begin{sol}
    Suppose $T' = 0$. Then for every $\varphi \in W'$ and every $v \in V$, we have
    \[(T'(\varphi))(v) = \varphi(T(v)) = 0.\]
    If $T(v) \neq 0$ for some $v \in V$, then $T(v) \in W$. Since $W$ is finite-dimensional, then by \autoref{ex:3f3}, there exists $\psi \in W'$ such that $\psi(T(v)) \neq 0$, which contradicts the fact that $(T'(\psi))(v) = \psi(T(v)) = 0$. Hence, $T(v) = 0$ for every $v \in V$, and so $T = 0$.

    Conversely, suppose $T = 0$. Then for every $\varphi \in W'$ and every $v \in V$, we have
    \[(T'(\varphi))(v) = \varphi(T(v)) = \varphi(0) = 0.\]
    Hence, $T' = 0$.
\end{sol}

\begin{exercise}[3f17]
    Suppose $V$ and $W$ are finite-dimensional and $T \in \lin{V, W}$. Prove that $T$ is invertible if and only if $T' \in \lin{W', V'}$ is invertible.
\end{exercise}

\begin{sol}
    Suppose $T$ is invertible. Then there exists $S \in \lin{W, V}$ such that $ST = I_V$ and $TS = I_W$. Taking the dual of both sides, we get $(ST)' = (I_V)'$ and $(TS)' = (I_W)'$. But $(ST)' = T'S'$ and $(TS)' = S'T'$, and $(I_V)' = I_{V'}$ and $(I_W)' = I_{W'}$. Hence, $T'S' = I_{V'}$ and $S'T' = I_{W'}$, which shows that $T'$ is invertible with inverse $S'$.

    Conversely, we prove the contrapositive. Suppose $T$ is not invertible. Then it is either not injective or not surjective.

    If $T$ is not injective, then there exists nonzero $v \in V$ such that $T(v) = 0$. By \autoref{ex:3f3}, there exists $\psi \in W'$ such that $\psi(v) = 1$. Since $\range T'$ contains all functionals of the form $T'(\varphi) = \varphi \circ T$ for $\varphi \in W'$, we have $(T'(\psi))(v) = \psi(T(v)) = \psi(0) = 0$, which contradicts $\psi(v) = 1$. Hence, $T'$ cannot be surjective, and therefore not invertible.

    If $T$ is not surjective, then there exists $w \in W$ such that $w \notin \range T$. Since $W$ is finite-dimensional, by \autoref{ex:3f4}, there exists $\varphi \in W'$ such that $\varphi(w) \neq 0$ and $\varphi$ vanishes on $\range T$. Then $T'(\varphi) = \varphi \circ T = 0$, which shows that $T'$ is not injective, and therefore not invertible. It follows that $T'$ is not invertible whenever $T$ is not invertible, which proves the contrapositive of the statement. Hence, if $T'$ is invertible, then $T$ must be invertible.
\end{sol}

\begin{exercise}[3f18]
    Suppose $V$ and $W$ are finite-dimensional. Prove that the map that takes $T \in \lin{V, W}$ to $T' \in \lin{W', V'}$ is an isomorphism of $\lin{V, W}$ onto $\lin{W', V'}$.
\end{exercise}

\begin{sol}
    Let $\Phi: \lin{V, W} \to \lin{W', V'}$ be the map defined by $\Phi(T) = T'$. We need to show that $\Phi$ is linear, injective, and surjective.

    Linearity: For $T_1, T_2 \in \lin{V, W}$ and $a \in \f$, we have
    \[\Phi(T_1 + T_2) = (T_1 + T_2)' = T_1' + T_2' = \Phi(T_1) + \Phi(T_2),\]
    and
    \[\Phi(aT_1) = (aT_1)' = aT_1' = a\Phi(T_1).\]
    Hence, $\Phi$ is linear.

    To see that $\Phi$ is injective, suppose $\Phi(T) = T' = 0$. By \autoref{ex:3f16}, this implies $T = 0$. Hence, $\Phi$ is injective.

    To see that $\Phi$ is surjective, note that since $V$ and $W$ are finite-dimensional, $\dim \lin{V, W} = \dim V \cdot \dim W = \dim W' \cdot \dim V' = \dim \lin{W', V'}$. A linear map between finite-dimensional vector spaces of the same dimension that is injective is also surjective. Hence, $\Phi$ is surjective and therefore an isomorphism.
\end{sol}

\begin{exercise}[3f19]
    Suppose $U \subseteq V$. Explain why
    \[U^0 = \set{\varphi \in V' \mid U \subseteq \nul \varphi}.\]
\end{exercise}

\begin{sol}
    By definition, the annihilator $U^0$ of $U$ is the set of all linear functionals $\varphi \in V'$ such that $\varphi(u) = 0$ for all $u \in U$. Equivalently, this means that $U \subseteq \nul \varphi$.
\end{sol}

\begin{exercise}[3f20]
    Suppose $V$ is finite-dimensional and $U$ is a subspace of $V$. Show that
    \[U = \set{v \in V \mid \varphi(v) = 0 \text{ for every } \varphi \in U^0}.\]
\end{exercise}

\begin{sol}
    Let $S = \set{v \in V \mid \varphi(v) = 0 \text{ for every } \varphi \in U^0}$. We need to show that $S = U$.

    First, if $v \in U$, then for any $\varphi \in U^0$, we have $\varphi(v) = 0$ by definition of the annihilator. Hence, $v \in S$, so $U \subseteq S$.

    Conversely, suppose $v \in S$. If $v \notin U$, then by \autoref{ex:3f4}, there exists a linear functional $\varphi \in V'$ such that $\varphi(u) = 0$ for all $u \in U$ and $\varphi(v) \neq 0$. But this $\varphi$ is in $U^0$, which contradicts $v \in S$. Hence, $v \in U$, and therefore $S \subseteq U$.

    Combining both inclusions, we get $S = U$, as desired.
\end{sol}

\begin{exercise}[3f21]
    Suppose $V$ is finite-dimensional and $U$ and $W$ are subspaces of $V$.
    \begin{subexercises}
        \item Prove that $W^0 \subseteq U^0$ if and only if $U \subseteq W$.
        \item Prove that $W^0 = U^0$ if and only if $U = W$.
    \end{subexercises}
\end{exercise}

\begin{sole}
    \item Suppose $W^0 \subseteq U^0$. If $u \in U$, then for any $\varphi \in W^0$, we have $\varphi(u) = 0$ because $\varphi \in U^0$. Hence, by \autoref{ex:3f20}, we have $u \in W$, so $U \subseteq W$. Conversely, if $U \subseteq W$ and $\varphi \in W^0$, then $\varphi(w) = 0$ for all $w \in W$, and in particular for all $u \in U$. Hence, $\varphi \in U^0$, so $W^0 \subseteq U^0$.
    \item Suppose $W^0 = U^0$. Then $W^0 \subseteq U^0$ and $U^0 \subseteq W^0$. By part (a), this implies that $U \subseteq W$ and $W \subseteq U$. Hence, $U = W$. Conversely, if $U = W$, then clearly $U^0 = W^0$.
\end{sole}

\begin{exercise}[3f22]
    Suppose $V$ is finite-dimensional and $U$ and $W$ are subspaces of $V$.
    \begin{subexercises}
        \item Show that $(U + W)^0 = U^0 \cap W^0$.
        \item Show that $(U \cap W)^0 = U^0 + W^0$.
    \end{subexercises}
\end{exercise}

\begin{sole}
    \item Let $\varphi \in (U + W)^0$. Then for any $u \in U$ and $w \in W$, we have $\varphi(u + w) = 0$. In particular, taking $w = 0$ gives $\varphi(u) = 0$ for all $u \in U$, so $\varphi \in U^0$. Similarly, taking $u = 0$ gives $\varphi(w) = 0$ for all $w \in W$, so $\varphi \in W^0$. Hence, $\varphi \in U^0 \cap W^0$, and we have shown that $(U + W)^0 \subseteq U^0 \cap W^0$.
    
    Conversely, let $\varphi \in U^0 \cap W^0$. Then for any $u \in U$ and $w \in W$, we have $\varphi(u) = 0$ and $\varphi(w) = 0$, so $\varphi(u + w) = \varphi(u) + \varphi(w) = 0$. Hence, $\varphi \in (U + W)^0$, and we have shown that $U^0 \cap W^0 \subseteq (U + W)^0$. Combining both inclusions, we get $(U + W)^0 = U^0 \cap W^0$.
    \item From part (a), we know that $(U + W)^0 = U^0 \cap W^0$. Then
    \[\dim(U + W)^0 = \dim V - \dim(U + W).\]
    Applying the dimension formula for a sum of subspaces, we have
    \[\dim(U^0 + W^0) = \dim U^0 + \dim W^0 - \dim(U^0 \cap W^0).\]
    Since $(U + W)^0 = U^0 \cap W^0$, we have $\dim(U^0 \cap W^0) = \dim(U + W)^0 = \dim V - \dim(U + W)$. Therefore,
    \[\dim(U^0 + W^0) = \dim U^0 + \dim W^0 - (\dim V - \dim(U + W)).\]
    Using the fact that $\dim U^0 = \dim V - \dim U$ and $\dim W^0 = \dim V - \dim W$, we get
    \begin{align*}
        \dim(U^0 + W^0) &= \dim U^0 + \dim W^0 - \dim(U^0 \cap W^0) \\
        &= (\dim V - \dim U) + (\dim V - \dim W) - (\dim V - \dim(U + W)) \\
        &= \dim V - \dim U + \dim V - \dim W - \dim V + \dim(U + W) \\
        &= \dim V - (\dim U + \dim W - \dim(U + W)) \\
        &= \dim V - \dim(U \cap W).
    \end{align*}
    Hence, $\dim(U^0 + W^0) = \dim(U \cap W)^0$. To see equality, we just need to prove one inclusion. If $\varphi = \varphi_U + \varphi_W \in U^0 + W^0$, then for any $v \in U \cap W$, we have $\varphi(v) = \varphi_U(v) + \varphi_W(v) = 0 + 0 = 0$, so $\varphi \in (U \cap W)^0$. Therefore, $U^0 + W^0 \subseteq (U \cap W)^0$. Since the dimensions are equal, we conclude that $(U \cap W)^0 = U^0 + W^0$.
\end{sole}

\begin{exercise}[3f23]
    Suppose $V$ is finite-dimensional and $\varphi_1, \ldots, \varphi_m \in V'$. Prove that the following three sets are equal to each other.
    \begin{subexercises}
        \item $\spn(\varphi_1, \ldots, \varphi_m)$
        \item $((\nul \varphi_1) \cap \cdots \cap (\nul \varphi_m))^0$
        \item $\set{\varphi \in V' \mid (\nul \varphi_1) \cap \cdots \cap (\nul \varphi_m) \subseteq \nul \varphi}$
    \end{subexercises}
\end{exercise}

\begin{sol}
    For ease of notation, let $A, B, C$ be the sets in the exercise, i.e.,
    \begin{align*}
        A &= \spn(\varphi_1, \ldots, \varphi_m),\\
        B &= ((\nul \varphi_1) \cap \cdots \cap (\nul \varphi_m))^0,\\
        C &= \set{\varphi \in V' \mid (\nul \varphi_1) \cap \cdots \cap (\nul \varphi_m) \subseteq \nul \varphi}.
    \end{align*}
    We will show that $A = B$ and $B = C$. Starting off with $B = C$, note that by definition,
    \begin{align*}
        B &= ((\nul \varphi_1) \cap \cdots \cap (\nul \varphi_m))^0 \\
          &= \set{\varphi \in V' \mid (\nul \varphi_1) \cap \cdots \cap (\nul \varphi_m) \subseteq \nul \varphi} = C,
    \end{align*}
    where the reformulation of the annihilator subspace was shown in \autoref{ex:3f19}.

    To show that $A = B$, we may use \autoref{ex:3f22} (b) repeatedly to rewrite $B$ as
    \[B = (\nul \varphi_1)^0 + \cdots + (\nul \varphi_m)^0.\]
    By \autoref{ex:3f6}, we have $(\nul \varphi_i)^0 = \spn(\varphi_i)$ for each $i = 1, \ldots, m$. Therefore,
    \[B = (\nul \varphi_1)^0 + \cdots + (\nul \varphi_m)^0 = \spn(\varphi_1) + \cdots + \spn(\varphi_m) = \spn(\varphi_1, \ldots, \varphi_m) = A.\]
    We conclude that $A = B$.
\end{sol}

\begin{exercise}[3f24]
    Suppose $V$ is finite-dimensional and $v_1, \ldots, v_m \in V$. Define a linear map $\Gamma: V' \to \f^m$ by $\Gamma(\varphi) = (\varphi(v_1), \ldots, \varphi(v_m))$.
    \begin{subexercises}
        \item Prove that $v_1, \ldots, v_m$ spans $V$ if and only if $\Gamma$ is injective.
        \item Prove that $v_1, \ldots, v_m$ is linearly independent if and only if $\Gamma$ is surjective.
    \end{subexercises}
\end{exercise}

\begin{sole}
    \item Suppose $v_1, \ldots, v_m$ spans $V$. If $\Gamma(\varphi) = 0$, then $\varphi(v_i) = 0$ for all $i = 1, \ldots, m$. Since $v_1, \ldots, v_m$ spans $V$, it follows that $\varphi(v) = 0$ for all $v \in V$, i.e., $\varphi = 0$. Hence, $\Gamma$ is injective.
    
    Conversely, suppose $\Gamma$ is injective. If $v_1, \ldots, v_m$ does not span $V$, then by \autoref{ex:3f4}, there exists a nonzero $\varphi \in V'$ that vanishes on $v_1, \ldots, v_m$. But then $\Gamma(\varphi) = 0$, contradicting the injectivity of $\Gamma$. Hence, $v_1, \ldots, v_m$ must span $V$.
    \item Suppose $v_1, \ldots, v_m$ is linearly independent. Extend this list to a basis $v_1, \ldots, v_m, v_{m + 1}, \ldots, v_n$ of $V$. Define linear functionals $\varphi_1, \ldots, \varphi_m \in V'$ by setting $\varphi_i(v_j) = \delta_{ij}$ for $i, j = 1, \ldots, m$ and $\varphi_i(v_j) = 0$ for $j = m + 1, \ldots, n$. Then for any $(a_1, \ldots, a_m) \in \f^m$, the linear functional $\varphi = a_1 \varphi_1 + \cdots + a_m \varphi_m$ satisfies $\Gamma(\varphi) = (a_1, \ldots, a_m)$. Hence, $\Gamma$ is surjective.
    
    Conversely, suppose $\Gamma$ is surjective. If $v_1, \ldots, v_m$ is not linearly independent, then there exist scalars $a_1, \ldots, a_m$, not all zero, such that $a_1 v_1 + \cdots + a_m v_m = 0$. Consider some index $j$ such that $a_j \neq 0$. Since $\Gamma$ is surjective, there exists $\varphi \in V'$ such that $\varphi(v_j) = 1$ and $\varphi(v_i) = 0$ for all $i \neq j$. Then
    \[0 = \varphi(a_1 v_1 + \cdots + a_m v_m) = a_1 \varphi(v_1) + \cdots + a_m \varphi(v_m) = a_j \neq 0,\]
    a contradiction. Hence, $v_1, \ldots, v_m$ must be linearly independent.
\end{sole}

\begin{exercise}[3f25]
    Suppose $V$ is finite-dimensional and $\varphi_1, \ldots, \varphi_m \in V'$. Define a linear map $\Gamma: V \to \f^m$ by $\Gamma(v) = (\varphi_1(v), \ldots, \varphi_m(v))$.
    \begin{subexercises}
        \item Prove that $\varphi_1, \ldots, \varphi_m$ spans $V'$ if and only if $\Gamma$ is injective.
        \item Prove that $\varphi_1, \ldots, \varphi_m$ is linearly independent if and only if $\Gamma$ is surjective.
    \end{subexercises}
\end{exercise}

\begin{sole}
    \item 
\end{sole}

\begin{exercise}[3f26]
    Suppose $V$ is finite-dimensional and $\Omega$ is a subspace of $V'$. Prove that
    \[\Omega = \set{v \in V \mid \varphi(v) = 0 \text{ for every } \varphi \in \Omega}^0.\]
\end{exercise}

\begin{exercise}[3f27]
    Suppose $T \in \lin{\pp_5(\r)}$ and $\nul T' = \spn(\varphi)$, where $\varphi$ is the linear functional on $\pp_5(\r)$ defined by $\varphi(p) = p(8)$. Prove that
    \[\range T = \set{p \in \pp_5(\r) \mid p(8) = 0}.\]
\end{exercise}

\begin{exercise}[3f28]
    Suppose $V$ is finite-dimensional and $\varphi_1, \ldots, \varphi_m$ is a linearly independent list in $V'$. Prove that
    \[\dim((\nul \varphi_1) \cap \cdots \cap (\nul \varphi_m)) = (\dim V) - m.\]
\end{exercise}

\begin{exercise}[3f29]
    Suppose $V$ and $W$ are finite-dimensional and $T \in \lin{V, W}$.
    \begin{subexercises}
        \item Prove that if $\varphi \in W'$ and $\nul T' = \spn(\varphi)$, then $\range T = \nul \varphi$.
        \item Prove that if $\psi \in V'$ and $\range T' = \spn(\psi)$, then $\nul T = \nul \psi$.
    \end{subexercises}
\end{exercise}

\begin{exercise}[3f30]
    Suppose $V$ is finite-dimensional and $\varphi_1, \ldots, \varphi_n$ is a basis of $V'$. Show that there exists a basis of $V$ whose dual basis is $\varphi_1, \ldots, \varphi_n$.
\end{exercise}

\begin{exercise}[3f31]
    Suppose $U$ is a subspace of $V$. Let $i: U \to V$ be the inclusion map defined by $i(u) = u$. Thus $i' \in \lin{V', U'}$.
    \begin{subexercises}
        \item Show that $\nul i' = U^0$.
        \item Prove that if $V$ is finite-dimensional, then $\range i' = U'$.
        \item Prove that if $V$ is finite-dimensional, then $i'$ is an isomorphism from $V' / U^0$ onto $U'$.
    \end{subexercises}
\end{exercise}

\begin{exercise}[3f32]
    The \textbf{double dual space} of $V$, denoted by $V''$, is defined to be the dual space of $V'$. In other words, $V'' = (V')'$. Define $\Lambda: V \to V''$ by
    \[(\Lambda v)(\varphi) = \varphi(v)\]
    for each $v \in V$ and each $\varphi \in V'$.
    \begin{subexercises}
        \item Show that $\Lambda$ is a linear map from $V$ to $V''$.
        \item Show that if $T \in \lin{V}$, then $T'' \circ \Lambda = \Lambda \circ T$, where $T'' = (T')'$.
        \item Show that if $V$ is finite-dimensional, then $\Lambda$ is an isomorphism from $V$ onto $V''$.
    \end{subexercises}
    \note{Suppose $V$ is finite-dimensional. Then $V$ and $V'$ are isomorphic, but finding an isomorphism from $V$ onto $V'$ generally requires choosing a basis of $V$. In contrast, the isomorphism $\Lambda$ from $V$ onto $V''$ does not require a choice of basis and thus is considered more natural.}
\end{exercise}

\begin{exercise}[3f33]
    Suppose $U$ is a subspace of $V$. Let $\pi: V \to V / U$ be the usual quotient map. Thus $\pi' \in \lin{(V / U)', V'}$.
    \begin{subexercises}
        \item Show that $\pi'$ is injective.
        \item Show that $\range \pi' = U^0$.
        \item Conclude that $\pi'$ is an isomorphism from $(V / U)'$ onto $U^0$.
    \end{subexercises}
\end{exercise}